%% P14 --- Logic, proof and reading theorems
%%
%% Every number in this program that the reader cannot do in their head comes
%% from code/p14_logic_proof.py and is pulled in with \val{}. The console
%% listing is a committed transcript written by that script.
%%
%% THE THESIS: a theorem is hypotheses plus a conclusion plus the quantifiers
%% binding them, and almost every misuse of one in this field is a reader
%% keeping the conclusion and dropping something else.
%%
%% WHAT THIS PROGRAM IS OWED, read out of the files rather than remembered:
%%   F10  gives and, or, not and De Morgan, and its filter example is
%%        continued here -- the counts fall out of the truth table and are
%%        gated against F10's committed ones.
%%   P13  DEFERS HERE BY NAME: it proves one direction of the topological
%%        order theorem and says "Program P14 is where reading and stating a
%%        theorem properly is taught".
%%   P05  used the union bound for its capacity figures and said in a
%%        rigourbox that how the bound was obtained matters. Section 5 uses
%%        the same bound on a different quantity.
%%   EVERY rigourbox in the book is this program's subject practised before
%%        it was named. NO TALLY IS STATED, anywhere: the count moves with
%%        every program written, and CLAUDE.md forbids a count of
%%        occurrences outright. Name the rule, not the number.
%%        And the rule is WIDER than "the proof is elsewhere", which was the
%%        first draft of it and is false: all fifteen were read, and while
%%        most point at a first course, P01, P02, P03 and P07 point at a later
%%        program and P04's second points at a MEASUREMENT nobody has made.
%%        What is true of every one: it names what the program is not doing
%%        and says where it is done instead.
%%
%% WHAT IT DOES NOT DO, and says so on the page: it does not train the reader
%% to WRITE proofs. The audience needs to read theorems, not prove them, and
%% pretending otherwise would be the third bad option after omitting rigour
%% and faking it.
%%
%% LESS TO COMPUTE THAN ANY OTHER PROGRAM, and that is a fact about the
%% subject. What can be computed is: every claim about implication and
%% quantifiers is settled by EXHAUSTIVE enumeration, which is a proof and not
%% a demonstration. What cannot be is labelled rather than dressed up.
%%
%% Twin: programs/pl/P14-logic-proof.tex. Frame for frame, macro for macro,
%% number for number; tools/parity.py compares them.

\program{Logic, proof and reading theorems}
\label{prog:P14}
\index{proof}
\index{proof!reading a theorem}

\begin{outcomes}
\outcome{1--7}{take a stated theorem apart into its hypotheses, its conclusion
  and the quantifiers binding them}
\outcome{8--14}{say why a statement and its contrapositive are one claim and
  its converse is another, and what a counterexample has to exhibit}
\outcome{15--19}{read the order of two quantifiers and say what changes when
  they are swapped}
\outcome{20--25}{recognise a proof by contradiction and a proof by induction,
  and say what would break each}
\outcome{26--31}{say, of three results this field misquotes, exactly which
  part is being dropped}
\end{outcomes}

\begin{quiz}
\item \emph{If it is a duck, it quacks.} Does it follow that a thing which
  quacks is a duck? \teachesat{9--10}
  \answerto{No. That is the \emph{converse} and it is a different claim; the
  two disagree in $\val{p14.differ}$ of the $\val{p14.rows}$ rows. What does
  follow is that a thing which does not quack is not a duck.}
\item A bound holds \enquote{with probability $\val{p14.conf}$ per cent}. You
  use it $\val{p14.uses}$ times. What is the chance all of them hold?
  \teachesat{28--29}
  \answerto{About $\val{p14.allhold}$ per cent, so it fails somewhere
  $\val{p14.anyfails}$ per cent of the time. A confidence quoted once is not a
  confidence in the conclusion you build out of twenty of them.}
\item \enquote{For every tolerance there is a network} and \enquote{there is a
  network for every tolerance}. Same claim? \teachesat{15--16}
  \answerto{No, and the second is far stronger. Swapping those two quantifiers
  changes what is allowed to depend on what, and universal approximation is
  the first of the two.}
\item A claim is checked on forty consecutive inputs and holds every time.
  What has been established? \teachesat{22--23}
  \answerto{That it holds on those forty. $n^{2}+n+41$ is prime for every $n$
  from $0$ to $\val{p14.euler.first}$ exclusive and composite at
  $\val{p14.euler.first}$, so forty confirmations can be worth nothing.}
\item A convergence theorem assumes the objective is convex. A paper cites it
  about a neural network. What has gone wrong? \teachesat{26--27}
  \answerto{A hypothesis was dropped. The conclusion is quoted and the
  condition it was proved under is not, and a neural network's loss is not
  convex \dash{} so the theorem simply does not apply.}
\end{quiz}

% ============================================================
\section{What a theorem is made of}
\label{sec:P14-parts}
\index{proof!hypotheses and conclusion}

\begin{fr}
This program has a smaller subject than it looks and a larger payoff. It will
not teach you to prove anything. It will teach you to read a theorem somebody
else has proved and say exactly what it does and does not claim, which is a
skill this field needs constantly and almost never teaches.

You have been watching it done. Every \emph{rigour} box in this book names
something the program is not doing and says where it is done instead \dash{}
usually a first course in the subject, sometimes a later program here, and once
a measurement nobody has made. There is one in this program too.

Look at that shape rather than at any one instance. What are the parts of a
theorem?
\dotline
\nextframe
\end{fr}

\begin{fr}
\begin{ansblock}
Its \textbf{hypotheses} \dash{} what is assumed. Its \textbf{conclusion}
\dash{} what is claimed. And the \textbf{quantifiers} \dash{} what each of
those ranges over.
\end{ansblock}

Three parts, and almost every misuse of a theorem in this field is one of them
going missing on the way to being quoted.

\begin{note}
It is worth being clear about which part a proof is. A proof is the argument
that the hypotheses force the conclusion; it is not one of the three parts and
you can use a theorem correctly without ever seeing it. That is the whole
reason this program is possible \dash{} and it is why a rigour box is not an
evasion. It withholds the argument and gives you all three parts, which is
everything you need in order to use the result and to notice when somebody
else is using it wrongly.
\end{note}

\mermaidfig{p14-parts-of-a-theorem}{The three parts, and which one gets
dropped.}{the three parts}
\end{fr}

\begin{fr}
Take a real one apart. Program~\ref{prog:P13} stated this:

\begin{quote}
A directed graph has a topological order if and only if it is acyclic.
\end{quote}

Name the hypotheses and the conclusion. There is a complication in the
sentence, and finding it is half the exercise.
\dotline
\nextframe
\end{fr}

\begin{fr}
\begin{ansblock}
It is \emph{two} theorems, and each is the other's converse.
\end{ansblock}

\enquote{If and only if} is an abbreviation for a pair of claims travelling
together: \emph{if it is acyclic then it has a topological order}, and
\emph{if it has a topological order then it is acyclic}. Each has its own
hypothesis, its own conclusion and, usually, its own proof \dash{} which is
exactly what Program~\ref{prog:P13} said when it proved one direction in a
sentence and handed the other here.

\textbf{The hypothesis in both is \enquote{directed}}, and it is doing real
work: an undirected graph with one edge has a walk out and back, so without
that word the theorem is false and the whole of Program~\ref{prog:P13}'s
section 4 evaporates. It is also the word a reader is likeliest to skip,
because it sits before the noun and reads like scenery.
\end{fr}

\begin{fr}
That is the pattern the rest of this program is about, so here it is by name.

\textbf{The commonest misuse of a theorem is to keep the conclusion and drop a
hypothesis.} Not to misunderstand the argument, not to doubt the result
\dash{} to quote what it concludes while leaving behind the conditions under
which it was proved.

Why would that be the common failure rather than some other one? Answer before
reading on; the reason is about people rather than about mathematics.
\nextframe
\end{fr}

\begin{fr}
\begin{ansblock}
Because the conclusion is the part anybody wanted. The hypotheses are what
somebody had to assume in order to get it, and nobody quotes a theorem for its
assumptions.
\end{ansblock}

A conclusion is quotable, transferable and short. A hypothesis is a
restriction, it is boring, and it is usually the technical clause in the middle
of the sentence. So the sentence that travels is the conclusion, and it travels
without the thing that made it true.

\begin{aibox}
This is not a hypothetical failure in this field, it is the normal one, and
section 5 works three instances. The shape is always the same: a result proved
under a condition arrives in a blog post, a slide or a review comment with the
condition gone, and it then gets applied to a system that does not satisfy it.

The defence is a habit rather than a technique. \textbf{When a result is
quoted at you, ask what it assumed} \dash{} and if the answer is not in the
sentence, the sentence is not yet a claim about your system.
\end{aibox}
\end{fr}

\begin{fr}
One thing this program deliberately does not do, stated plainly because the
alternative is to leave you guessing.

\textbf{It does not teach you to write proofs.} Not a shortened version, not
the easy ones. Writing proofs is a craft that takes a course and a year, and
almost nobody doing this job needs it.

\begin{rigourbox}
The three options for a book like this were: leave rigour out entirely, which
is what most practical mathematics for engineers does and which leaves the
reader unable to read the literature; teach proof-writing, which is a
different book for a different reader; or teach \emph{reading}, which is the
skill actually needed and is a fraction of the work.

This is the third, and the boundary is worth stating so that you know what you
will and will not have. After this program you will be able to say what a
theorem claims, what it assumes, what would refute it and what it is silent
about. You will not be able to prove it, and you will not need to.
\end{rigourbox}
\end{fr}

% ============================================================
\section{Implication is not equivalence}
\label{sec:P14-implication}
\index{proof!implication}

\begin{fr}
Almost every hypothesis-and-conclusion pair is an \emph{implication}: if $P$,
then $Q$. Program~\ref{prog:F10} gave \emph{and}, \emph{or} and \emph{not};
this is the fourth connective and it is the one people get wrong.

There are $\val{p14.rows}$ cases, since $P$ and $Q$ are each true or false.
In how many of them is \enquote{if $P$ then $Q$} \emph{false}?
\dotline
\nextframe
\end{fr}

\begin{fr}
\ans{One}
Only when $P$ holds and $Q$ fails. If $P$ does not hold the implication makes
no claim at all and is counted true, which reads as a technicality and is the
whole of what the connective means.

\begin{center}
\begin{tabular}{cccc}
\toprule
$P$ & $Q$ & $P \Rightarrow Q$ & $Q \Rightarrow P$ \\
\midrule
F & F & T & T \\
F & T & T & F \\
T & F & F & T \\
T & T & T & T \\
\bottomrule
\end{tabular}
\end{center}

\textbf{So \enquote{if $P$ then $Q$} is exactly the statement that there is no
case with $P$ holding and $Q$ failing.} Nothing about cause, nothing about
what happens when $P$ is false, nothing about $Q$ being interesting.

Now read the fourth column, which is the converse $Q \Rightarrow P$. In how
many rows does it disagree with the third?
\nextframe
\end{fr}

\begin{fr}
\ans{$\val{p14.differ}$ of the $\val{p14.rows}$}
So they are different claims, and knowing one tells you nothing about the
other.

\begin{trapbox}
\textbf{Affirming the consequent} is taking $P \Rightarrow Q$, observing $Q$,
and concluding $P$. It is the second row of the table, where the implication
holds and the converse does not.

In this field it usually arrives as a diagnosis. \emph{Overfitting drives the
training loss down, and our training loss is down, so we are overfitting.} The
first clause may be perfectly true and the conclusion still does not follow:
a model that has genuinely learnt the task also has a low training loss.

The test is mechanical and takes a second. Write the claim as
\emph{if $P$ then $Q$}, write down what was observed, and check which of the
two letters it is. If it is $Q$, you have the converse and you have nothing.
\end{trapbox}

There is a mirror-image error, and it is the one that produces false
reassurance rather than a false diagnosis.

Take $P \Rightarrow Q$ again and suppose you observe that $P$ is
\emph{false}. What follows about $Q$?
\nextframe
\end{fr}

\begin{fr}
\ans{Nothing at all}
Rows one and two of the table: with $P$ false the implication holds either
way, so $Q$ is free. Concluding $\lnot Q$ from $\lnot P$ is \textbf{denying
the antecedent}, and it is the converse's error seen from the other side
\dash{} indeed $\lnot P \Rightarrow \lnot Q$ \emph{is} the converse's
contrapositive, so the two mistakes are one mistake.

It reads as safety rather than as a claim, which is what makes it slip
through. \emph{A model that has seen the test set will score suspiciously
well. This one has not seen it, so the score is trustworthy.} The second
sentence removes one explanation for a high score and leaves every other one
untouched.

Now the third form. Does \enquote{if not $Q$ then not $P$} agree with
\enquote{if $P$ then $Q$}?
\dotline
\nextframe
\end{fr}

\begin{fr}
\ans{Yes, in all $\val{p14.rows}$ rows}
That is the \textbf{contrapositive}, and it is not a related statement but the
same one written differently \dash{} which the script checks by enumerating
every row, so it is a proof rather than an illustration.

It is worth having because it is often the easier half to establish.
\emph{Every duck quacks} and \emph{nothing that fails to quack is a duck} are
one claim; if the second is easier to check, checking it has proved the first.
That is what a proof by contraposition is, and it is not a trick.

\mermaidfig{p14-two-rewritings}{One of the two rewritings is the
same claim and one is a different one.}{the two rewritings}
\end{fr}

\begin{fr}
One more thing falls out of the table, and it is the most practical thing in
the section. What would it take to \emph{refute} \enquote{if $P$ then $Q$}?
\nextframe
\end{fr}

\begin{fr}
\begin{ansblock}
One case with $P$ holding and $Q$ failing. Exactly one, and nothing else is
required.
\end{ansblock}

Because the negation of $P \Rightarrow Q$ is $P \wedge \lnot Q$ \dash{} which
is Program~\ref{prog:F10}'s De Morgan one step on, and the script checks it
across all $\val{p14.rows}$ rows.

\textbf{That is what a counterexample is}, and the asymmetry is worth
noticing: establishing an implication may need an argument about every case,
and destroying one needs a single exhibit. It is why a reviewer who produces
one instance has finished the discussion, and why \enquote{but it usually
works} is not a reply.

\begin{note}
Program~\ref{prog:F10} measured this same connective going wrong in a filter:
of $\val{p14.filter.total}$ records, the right condition keeps
$\val{p14.filter.right}$ and \code{not spam and not short} keeps
$\val{p14.filter.wrong}$, losing $\val{p14.filter.lost}$ with no symptom. This
program's script rebuilds that population from the truth table and fails the
build if either count disagrees with the one that program committed \dash{} so
the logic here and the measurement there cannot come apart.
\end{note}
\end{fr}

% ============================================================
\section{Quantifiers, and the order that changes the claim}
\label{sec:P14-quantifiers}
\index{proof!quantifiers}

\begin{fr}
The second part of a theorem people drop is the quantifier, and specifically
its \emph{order}.

Two sentences about a relation between $x$ and $y$:

\begin{itemize}
\item \textbf{For every $x$ there is a $y$} such that $R(x,y)$.
\item \textbf{There is a $y$} such that for every $x$, $R(x,y)$.
\end{itemize}

They use the same five words. Are they the same claim? If not, which is
stronger, and why?
\dotline
\nextframe
\end{fr}

\begin{fr}
\begin{ansblock}
Different, and the second is stronger. In the first the $y$ may depend on the
$x$; in the second one $y$ has to work for all of them at once.
\end{ansblock}

The script settles it by exhibiting a relation on $\val{p14.dom}$ elements
\dash{} each $x$ related to the next one round a cycle, so
$\val{p14.rel}$ of the $\val{p14.pairs}$ pairs \dash{} for which the first
sentence is true and the second is false. One witness is a proof that the two
are different claims, which is section 2's asymmetry doing work.

\begin{note}
Only one direction is safe, and it is worth knowing which. \textbf{There-is-a-
$y$-for-all-$x$ always implies for-all-$x$-there-is-a-$y$}, never the reverse:
the universal witness serves as each individual one. The script checks that
over \emph{every} relation on a three-element set \dash{} all
$\val{p14.qchecked}$ of them \dash{} which is again an exhaustive proof rather
than a sample.
\end{note}
\end{fr}

\begin{fr}
Now the one this field quotes most. The universal approximation theorem, in
the shape it is usually stated:

\begin{quote}
For every continuous function $f$ on a bounded region and every tolerance
$\varepsilon > 0$, there exists a network with one hidden layer whose output
is within $\varepsilon$ of $f$ everywhere on that region.
\end{quote}

Which order is that, and what is therefore allowed to depend on what?
\dotline
\nextframe
\end{fr}

\begin{fr}
\begin{ansblock}
For-all, then there-exists. So the network is allowed to depend on both $f$
and $\varepsilon$ \dash{} and does.
\end{ansblock}

Read that again with section 2's habit: what does the sentence actually claim?
That for each demand, \emph{something} meets it. And then read what is not in
the sentence at all.

\begin{itemize}
\item How large the network has to be.
\item How that size grows as $\varepsilon$ shrinks, or as the region gains
  dimensions.
\item How the weights are to be found.
\item Whether any training procedure finds them.
\item How much data that would take.
\end{itemize}

$\val{p14.silent}$ things, and every one of them is what somebody quoting the
theorem was actually asking about.

\mermaidfig{p14-exists-is-not-found}{What an existence theorem promises, and
the three questions it does not touch.}{exists is not found}
\end{fr}

\begin{fr}
\begin{trapbox}
\textbf{\enquote{Neural networks can approximate any function} is the
conclusion with the quantifier and the silence both removed.} The theorem is
true, the sentence is a fair summary of its conclusion, and using it to
justify a choice of architecture is a mistake of a kind this program can name
precisely: it treats an existence claim as a construction.

The parallel that makes it obvious is that a polynomial of high enough degree
also approximates any continuous function on a bounded region, and nobody
concludes from that that polynomials are a good model class. What decides the
question is everything the theorem is silent about.
\end{trapbox}

\begin{warning}
\textbf{How the required size grows is a real question with real answers, and
this book has neither run the experiment nor surveyed that literature.} There
are results in both directions \dash{} constructions that are efficient for
some function classes, and lower bounds that are exponential in the dimension
for others.

Saying which applies to any particular architecture is beyond what this book
has done, so it is not said. What is established here is the shape of the
statement, which is a fact about the quantifiers and needs no experiment.
\end{warning}
\end{fr}

% ============================================================
\section{Three shapes, and what each one would take to break}
\label{sec:P14-shapes}
\index{proof!induction}
\index{proof!contradiction}

\begin{fr}
You do not need to write proofs and you do need to recognise three shapes,
because each one tells you what a refutation would have to look like.

A \textbf{direct} proof assumes the hypotheses and reaches the conclusion.
A \textbf{proof by contradiction} assumes the hypotheses and the negation of
the conclusion and derives an impossibility. A \textbf{proof by induction}
establishes a base case and a step from each case to the next.

The second one is section 2's result being used. Assuming the negation of
\enquote{if $P$ then $Q$} means assuming what, exactly?
\nextframe
\end{fr}

\begin{fr}
\ans{$P$ and not $Q$}
Which is why a proof by contradiction of an implication always begins by
supposing the hypothesis holds and the conclusion fails \dash{} it is not a
stylistic choice, it is the only thing the negation can be.

Program~\ref{prog:P13} used one without labelling it: if a directed graph had
a cycle, each vertex on it would have to come after the next one round, and a
finite list cannot do that. Assume the hypothesis, assume the conclusion
fails, reach something impossible.
\end{fr}

\begin{fr}
Induction is the one that matters most here, because it is the one a habit
from engineering makes people think they have.

A claim about every natural number. You check it for $n = 0$, for $n = 1$, and
onward, and it holds every time you look. Why is that not a proof, and what
would make it one?
\dotline
\nextframe
\end{fr}

\begin{fr}
\begin{ansblock}
Because you checked finitely many and the claim is about infinitely many. What
makes it a proof is a \emph{step}: an argument that if it holds at $n$ it
holds at $n+1$.
\end{ansblock}

The base case and the step together reach every $n$, and neither alone reaches
any. That is the whole shape, and recognising it tells you what to attack: a
failed induction is almost always a step that quietly assumes something
extra, not a wrong base case.

Here is why the distinction is not academic. Consider $n^{2} + n + 41$.

\transcript{p14-forty-confirmations}

\begin{trapbox}
\textbf{It is prime for every $n$ from $0$ to $\val{p14.euler.first}$
exclusive, and composite at $\val{p14.euler.first}$}, where it is
$\val{p14.euler.at}$ \dash{} which is $\val{p14.euler.factor}$ squared.

Forty consecutive confirmations. Any test suite would have passed it; any
reviewer would have believed it; and the claim is false.

The honest half of this is worth as much as the trap. In engineering a
thousand tests is often exactly the right instrument, and this book has used
enumeration as a proof several times \dash{} but only over domains that are
\emph{finite and exhausted}: four rows, $\val{p14.qchecked}$ relations,
$\val{p14.pairs}$ pairs. \textbf{Checking every case is a proof; checking many
cases is evidence}, and the difference is not how many, it is whether any were
left.
\end{trapbox}
\end{fr}

\begin{fr}
Each shape fails in its own way, and knowing which way is what makes
recognising the shape worth anything. Take them in turn: if a proof of one of
these three kinds is wrong, where is the wrongness most likely to be?
\dotline
\nextframe
\end{fr}

\begin{fr}
\begin{ansblock}
In a \emph{direct} proof, a step that does not follow. In one by
\emph{contradiction}, the negation being taken wrongly. In one by
\emph{induction}, the step \dash{} almost never the base case.
\end{ansblock}

\begin{itemize}
\item \textbf{Direct}: some line does not follow from the ones above it, and
  the usual culprit is a step that silently uses a property nobody assumed.
\item \textbf{Contradiction}: the negation is wrong. Negating
  \emph{for every $x$ there is a $y$} gives \emph{there is an $x$ with no
  $y$}, and getting that wrong makes the whole argument prove something else.
  Section 3's quantifiers and section 2's $P \wedge \lnot Q$ are what stop
  it happening.
\item \textbf{Induction}: the step holds only for $n$ large enough, or it
  quietly needs the claim for \emph{all} smaller $n$ when only the previous
  one is available. The base case is nearly always fine, because it is the
  part that gets checked.
\end{itemize}

\textbf{So the three shapes are three questions to ask}, and none of them
requires you to follow the argument in detail. What was negated? Does the step
depend on $n$ being large? Which line uses something that was not assumed?
That is what recognising the shape buys, and it is the whole of what this
section claimed it would.
\end{fr}

% ============================================================
\section{Three results this field misquotes}
\label{sec:P14-misquoted}
\index{proof!misquoted results}

\begin{fr}
Everything so far has been machinery. Here it is used, on three statements you
will meet.

The first. A convergence theorem says that gradient descent with a suitable
step size converges to the global minimum, \emph{provided the objective is
convex}. A paper cites it about a neural network.

Say what has gone wrong, and be precise about which of the three parts is
involved.
\dotline
\nextframe
\end{fr}

\begin{fr}
\begin{ansblock}
A hypothesis has been dropped. The conclusion is quoted; the condition it was
proved under is not; and a neural network's loss is not convex, so the theorem
does not apply at all.
\end{ansblock}

Not \enquote{applies approximately} and not \enquote{applies in spirit}. An
implication whose hypothesis fails makes no claim \dash{} that is the first
two rows of section 2's table, and it is the whole content of them.

\begin{note}
What survives is worth being fair about, because \enquote{the theorem does not
apply} is sometimes taken to mean \enquote{the method does not work}, and that
does not follow either. Gradient descent demonstrably works on non-convex
losses. What is gone is the \emph{guarantee}, and with it the right to say
\enquote{it will converge to the minimum} rather than \enquote{it has, on
every problem we have tried}.

That is a smaller claim and it is an honest one, and this book's own rule
\dash{} a claim is measured, attributed, or labelled as judgement \dash{} is
the same discipline applied to a sentence rather than to a theorem.
\end{note}
\end{fr}

\begin{fr}
The second, and it is the one that costs money. A generalisation bound holds
\enquote{with probability at least $\val{p14.conf}$ per cent over the draw of
the training data}.

You are reading a paper that applies such a bound $\val{p14.uses}$ times
\dash{} once per experiment, or once per model in a comparison. Assuming the
draws are independent, what is the chance that all $\val{p14.uses}$ hold?
\dotline
\nextframe
\end{fr}

\begin{fr}
\ans{About $\val{p14.allhold}$ per cent}
So the chance that at least one of them fails is $\val{p14.anyfails}$ per
cent. \textbf{More likely than not, one of the twenty statements in that
table is simply false}, and nothing in the paper marks which.

Two more numbers make the shape clear.

\begin{center}
\begin{tabular}{ll}
\toprule
uses before it is a coin flip & $\val{p14.flip}$ \\
per-use confidence needed for $95\%$ across $\val{p14.uses}$
  & $\val{p14.needed}$ per cent \\
\bottomrule
\end{tabular}
\end{center}

\textbf{A $\val{p14.conf}$ per cent bound is a coin flip after
$\val{p14.flip}$ uses}, and to have twenty of them hold together you would
need each to hold at $\val{p14.needed}$ per cent \dash{} a failure probability
of $\val{p14.neededdelta}$ per cent rather than $5$, nearly twenty times
tighter.

\begin{note}
The independence assumption is doing work above, and it is usually false: the
twenty experiments often share data. The bound that needs no independence is
the union bound, which Program~\ref{prog:P05} used for its capacity figures
\dash{} and here it gives $\val{p14.uses} \times 5 = \val{p14.union}$ per
cent, which is to say nothing at all.

That is not a failure of the union bound; it is the honest statement that with
twenty uses at this confidence, no guarantee survives. Both calculations agree
about what to do, which is to tighten the per-use bound rather than to hope.
\end{note}
\end{fr}

\begin{fr}
The third is universal approximation, and section 3 has already taken it
apart, so what is left is the pattern the three share.

\begin{aibox}
\textbf{Convexity} is a dropped hypothesis. \textbf{With high probability} is
a dropped quantifier \dash{} the bound is quantified over draws of the data,
and quoting it about the dataset in hand silently swaps a statement about a
distribution for one about a sample. \textbf{Universal approximation} is a
dropped silence: the theorem says a thing exists and the quotation implies it
can be found.

One of each: a hypothesis, a quantifier, and the gap between existence and
construction. They are the three parts of section 1, and there is not a fourth
kind of mistake here because there is not a fourth part.
\end{aibox}

So the reading procedure is short enough to use on a train. What is assumed?
What is claimed? What does each range over? And what is the statement silent
about that the person quoting it needed?
\end{fr}

\begin{fr}
One closing observation, which is what this program has really been about.

The rigour boxes in this book were never apologies. Every one of them names
what is being asserted, admits the rest of it is elsewhere, and says where
\dash{} which is precisely the discipline this program has now made explicit,
practised from the Foundation part onward before it was named.

\textbf{You are not short of proof; you are short of the habit of asking what
a result assumed.} That habit costs a sentence, it applies to a paper, a
colleague and a documentation page alike, and it is the only part of this
subject you will use every week.
\end{fr}

% ============================================================
\begin{summarybox}
\sumitem{1--2}{\textbf{A theorem is hypotheses, a conclusion and the
  quantifiers binding them.} A proof is none of the three, which is why a
  result can be used correctly without ever seeing one \dash{} and why a
  rigour box is not an evasion.}
\sumitem{3--4}{\enquote{If and only if} is \textbf{two} theorems travelling
  together, each the other's converse. In Program~\ref{prog:P13}'s
  topological-order theorem the load-bearing hypothesis is
  \emph{directed}, and it is the word a reader skips.}
\sumitem{5--6}{\textbf{The commonest misuse is to keep the conclusion and drop
  a hypothesis}, because the conclusion is the part anybody wanted. When a
  result is quoted at you, ask what it assumed.}
\sumitem{7}{This program does not teach proof-writing, deliberately: the
  three options were to omit rigour, to teach writing, or to teach
  \emph{reading}, and reading is the skill this audience needs.}
\sumitem{8--9}{\textbf{$P \Rightarrow Q$ is false in exactly one of the
  $\val{p14.rows}$ rows} \dash{} $P$ true and $Q$ false. It says nothing about
  cause and nothing about the case where $P$ fails.}
\sumitem{9--10}{\textbf{The converse is a different claim}, disagreeing in
  $\val{p14.differ}$ rows. Observing $Q$ and concluding $P$ is affirming the
  consequent: \enquote{the training loss is down, so we are overfitting}.}
\sumitem{10--11}{\textbf{And $\lnot P$ tells you nothing about $Q$.}
  Concluding $\lnot Q$ from $\lnot P$ is denying the antecedent, which is the
  converse's error from the other side \dash{} it reads as reassurance, which
  is what lets it through.}
\sumitem{11--12}{\textbf{The contrapositive is the same claim}, in all
  $\val{p14.rows}$ rows, which is why proving it proves the original.}
\sumitem{13--14}{\textbf{A counterexample is one case with the hypothesis
  holding and the conclusion failing}, because
  $\lnot(P \Rightarrow Q) = P \wedge \lnot Q$. Establishing may need every
  case; refuting needs one.}
\sumitem{15--16}{\textbf{Quantifier order changes the claim.} There-is-a-$y$-
  for-all-$x$ is the stronger, and implies the other; the reverse fails, as a
  relation on $\val{p14.dom}$ elements shows and all
  $\val{p14.qchecked}$ relations on three confirm.}
\sumitem{17--18}{\textbf{Universal approximation is for-all-then-there-exists},
  so the network may depend on the function and the tolerance. It is silent on
  $\val{p14.silent}$ things, including how large the network must be and
  whether training finds it.}
\sumitem{19}{\enquote{Networks can approximate any function} treats an
  existence claim as a construction. A high-degree polynomial does the same
  and nobody argues from it.}
\sumitem{20--21}{Three shapes to recognise: direct, contradiction, induction.
  \textbf{The negation of an implication is $P \wedge \lnot Q$}, which is why
  a proof by contradiction of one always starts there.}
\sumitem{22--23}{\textbf{Checking every case is a proof; checking many cases is
  evidence.} $n^{2}+n+41$ is prime for $\val{p14.euler.first}$ consecutive
  values and composite at $\val{p14.euler.first}$, where it is
  $\val{p14.euler.factor}$ squared.}
\sumitem{24--25}{\textbf{Each shape fails in its own place}: a direct proof
  at a step using something nobody assumed, a contradiction at the negation, an
  induction at the step and almost never at the base case. Three questions, and
  none needs the argument followed in detail.}
\sumitem{26--27}{\textbf{A convergence theorem assuming convexity says nothing
  about a neural network.} What is lost is the guarantee, not the method
  \dash{} which is this book's own rule about measurement applied to a
  sentence.}
\sumitem{28--29}{\textbf{A $\val{p14.conf}$ per cent bound used
  $\val{p14.uses}$ times holds throughout only $\val{p14.allhold}$ per cent of
  the time}, is a coin flip after $\val{p14.flip}$ uses, and would need
  $\val{p14.needed}$ per cent each to give $95\%$ across twenty.}
\sumitem{30--31}{The three misquotations are one of each part: a dropped
  hypothesis, a dropped quantifier, and the gap between existence and
  construction. \textbf{There is no fourth kind of mistake because there is no
  fourth part.}}
\end{summarybox}

\canyou

\begin{testexercises}
\item \emph{If a model is memorising, its validation loss rises.} The
  validation loss has risen. What follows?
  \answerto{Nothing about memorising. That is the converse, and it disagrees
  with the original in $\val{p14.differ}$ of the $\val{p14.rows}$ rows \dash{}
  a distribution shift or a learning rate that is too high will raise it too.
  What does follow, from the contrapositive, is that a model whose validation
  loss has not risen is not memorising \dash{} which is a real and useful
  deduction from the same sentence.}
\item Write the negation of \emph{every batch in this dataset contains at
  least one positive example}, and say what you would have to produce to
  refute the original.
  \answerto{\emph{Some batch contains no positive example.} Negating a
  for-all gives a there-exists and flips what is inside. To refute the
  original you produce one such batch \dash{} not a statistic about batches,
  one batch.}
\item A bound holds with probability $99\%$ per use. Roughly how many uses
  before the chance all hold drops below a half?
  \answerto{About $69$: $\num{0.99}^{n} < \num{0.5}$ needs
  $n > \ln \num{0.5} / \ln \num{0.99}$. Compare $\val{p14.flip}$ uses at
  $\val{p14.conf}$ per cent \dash{} tightening the per-use bound from $5\%$ to
  $1\%$ buys about five times as many uses, which is the trade the paper's
  author had to make and usually does not state.}
\item \emph{There is a learning rate that works for every architecture}
  against \emph{for every architecture there is a learning rate that works}.
  Which would you rather have, and which is plausible?
  \answerto{The first is far stronger and is what a default in a library
  claims implicitly; the second is nearly vacuous and is what anybody would
  actually defend. The whole difference is the order of two quantifiers, and
  a tuning guide that slides from the second to the first is making a much
  bigger claim than it can support.}
\item Somebody says a property was verified on ten thousand random inputs.
  What have they established, and when is that enough?
  \answerto{That it holds on those ten thousand. It is enough when the domain
  is finite and they were all of it \dash{} which is how this program proves
  its statements about $\val{p14.rows}$ rows and
  $\val{p14.qchecked}$ relations \dash{} and it is evidence rather than proof
  otherwise. $n^{2}+n+41$ survives forty consecutive checks and is false.}
\end{testexercises}

\begin{furtherproblems}
\item Show from the truth table that $P \Rightarrow Q$ and
  $\lnot P \vee Q$ are the same statement.
  \answerto{Both are false in exactly the row with $P$ true and $Q$ false, and
  true in the other three. This is why the negation is $P \wedge \lnot Q$: De
  Morgan applied to $\lnot P \vee Q$, which is Program~\ref{prog:F10}'s rule
  doing the work.}
\item Why is \emph{if $P$ then $Q$} counted true whenever $P$ is false, when
  that feels like it should be undefined?
  \answerto{Because the alternative breaks the connective. \emph{Every element
  of the empty set is prime} has to be true for \emph{every element of $A$ is
  prime and every element of $B$ is prime implies every element of
  $A \cup B$ is prime} to hold in all cases. The convention is chosen so that
  the rules compose, which is the same reason the empty product is $1$ in
  Program~\ref{prog:F04}.}
\item A theorem's hypotheses are $A$ and $B$; its conclusion is $C$. Your
  system satisfies $A$ but not $B$. What, if anything, may you conclude?
  \answerto{Nothing from this theorem, in either direction. It does not say
  $C$ holds and it does not say $C$ fails; the hypotheses failed, so the
  implication makes no claim. Concluding that $C$ fails is denying the
  antecedent, which is the converse's error wearing a different hat.}
\item Universal approximation is stated for one hidden layer. What would a
  version stated for two hidden layers have to add to be a stronger result?
  \answerto{Something about size or about findability. Purely as an existence
  claim it would be weaker if anything \dash{} anything one layer can do, two
  can \dash{} so a theorem about depth is only interesting if it bounds how
  large the network must be, or how the size grows. That is the question the
  one-layer statement is silent about, and it is where the real literature is.}
\item Find three rigour boxes in programs you have already read. Pick the
  shape they share and say what a reader is entitled to do with a result
  presented that way.
  \answerto{Each gives the statement, admits that the rest of it is elsewhere
  \dash{} a proof, a later program, or in one case a measurement nobody has
  made \dash{} and names where. A reader may use the result as if they had seen
  the proof
  \dash{} the three parts are all present \dash{} and may not extend it,
  weaken a hypothesis or apply it outside its quantifiers, because those are
  the moves the argument would have had to justify.}
\end{furtherproblems}
